\documentclass[12pt]{article}
\usepackage[margin=1in]{geometry}      % 四边各1英寸边距
\usepackage{setspace}
\doublespacing                         % 两倍行距
\usepackage{indentfirst}               % 每个段落（包括节后首段）首行缩进
\usepackage{times}
%\usepackage{txfonts}
%\usepackage{newtxtext, newtxmath}      
%\usepackage{newtxtext}
\usepackage{graphicx}
\usepackage{amsmath}
\usepackage{amssymb}
\usepackage{physics}
%\usepackage{tikz-feynman}
\usepackage{natbib}
\usepackage{hyperref}
\usepackage{amsfonts}



\title{Notes}
\author{LKC-HZ}
%\date{April 2026}
\date{}

\begin{document}
	
	\citestyle{super}
	\maketitle
	
	\tableofcontents
	
	
	\section{What is $SU(3)$ Group?}
	\subsection{Group Theory \label{group theory basics section}}
	A group is a set of elements $\{g_i\}$ together with a multiplication rule $g_i \times g_j = g_k$ that tells how each pair of elements is combined to give a third.
	
	A representation of a group $G$ is its realization in a linear space, where each of the group elements $g$ corresponds to a $d \times d$ matrix $D(g)$ and the composition of the group elements is represented by matrix multiplication: 
	\begin{align}
		D(g_1)D(g_2) = D(g_1g_2).    
	\end{align}
	For example, the three‑dimensional rotation matrices are representations of the SO(3) group elements. 
	
	Two representations $D$ and $D'$ are equivalent if they can be related by a similarity transformation:
	\begin{equation}
		D' = S^{-1}DS.
	\end{equation}
	The two equivalent matrices represent the very same thing, but they are of different basis. If we take the trace of $D'$:
	\begin{align}
		\tr(D') = \tr(S^{-1}DS) \rightarrow \tr(AB) = \tr(BA) \rightarrow \tr(D') = \tr(DS^{-1}S ) = \tr(D).
	\end{align}
	In group theory, the trace of a representation matrix is called character $\chi(g)$:
	\begin{equation}
		\chi(g) = \tr(D(g)).
	\end{equation}
	The character is basis-independent. 
	
	
	Given two equivalent representations $D^{(r)}$ and $D^{(s)}$ of a finite group $G$, which is $r$-dimensional and $s$-dimensional representation respectively, we have (\cite{azee2016group}, p. 103)
	\begin{equation}
		\sum_g D^{(r)\dagger}(g)^i_{\ j}D^{(s)}(g)^k_{\ l} = \frac{N(G)}{r}\delta^{rs}\delta ^i_{\ l}\delta^k_{\ j}. \label{GOT}
	\end{equation}
	Here, the Kronecker delta is meaningful only if $r = s$. It is also important to note that here we are summing over all $g \in G$; no contractions are made between $r, s, i, j, k$ or $l$. $\delta^i_{\ l}$ and $\delta^k_{\ j}$ imply that the summatiom has a non-zero result only if the two matrix elements taken from $D^{(r)}$ and $D^{(s)}$ stays in the same position; since $D^{(r)}$ is taken adjoint, $i=l$ and $j=k$ instead of $i=k$ and $j = l$. If $G$ in Eq. (\ref{GOT}), say, is a 2-dimensional rotational group ($SO(2)$ group, which is not finite but compact), its representation is rotational matrices. Taken a specific $g$, the product $D(g)^i_{\ j}D(g)^k_{\ l}$ depends on parameters of $g$, namely, the rotation angle. However, when we take the sum of $D(g)$'s of all possible $g$'s, angles are averaged so that the result is orientation-blind. Eq. (\ref{GOT}) can also be written in terms of character by taking $i = j$ and $k = l$ (\cite{azee2016group}, p. 104):
	\begin{align}
		\sum_cn_c(\chi^{(r)}(c))^*\chi^{(s)}(c) = N(G)\delta^{rs}, \label{GOT character version}
	\end{align}
	where $n_c$ is the number of elements belonging to class $c$. $n_c$ comes from the fact that the character of every individual representation in a class is the same as others. And since now we are summing over classes, the character each class should be multiplied by $n_c$ so that the result matches Eq. (\ref{GOT}) in which we summed over all group elements.
	
	A representation $D$ is reducible (irreducible otherwise) if it can be block-diagonalized by a similariy transformation:
	\begin{align}
		S^{-1}DS = \begin{pmatrix}
			D_1 &  & \\
			& D_2 & \\
			&&\dots
		\end{pmatrix} = D_1 \oplus D_2 \oplus \dots 
	\end{align}
	The other way round, we can also use the direct sum to create a larger equivalent representation matrix from smaller matrices. If $D(g)$ is an irreducible representation of a finite group $G$ and if there is some matrix $A$ such that $AD(g) = D(g)A$ for all $g$, then $A = \lambda I$ for some constant $\lambda$ (Schur's Lemma). Suppose now we have a class $c$ of reducible representations, the character of the class $\chi(c)$ can be written in terms of a summation of the character of irreducible representations $\chi^{(r)}(c)$ (this is reasonable since the trace of a block-diagonalized matrix equals that of every individual block adding together):
	\begin{equation}
		\chi(c) = \sum_r n_r\chi^{(r)}(c).
	\end{equation}
	Substitute into Eq. (\ref{GOT character version}) gives
	\begin{align}
		\sum_cn_c\chi^*(c)\chi(c) = N(G)\sum_r(n_r)^2 \\
		\Rightarrow \sum_r(n_r)^2 = \frac{1}{N(G)}\sum_c n_c\chi^*(c)\chi(c). \label{irreducible = 1 equation}
	\end{align}
	Since $n_r$ is the number of irreducible representations $\chi^{(r)}(c)$ that compose the reducible $\chi(c)$, the LHS of Eq. (\ref{irreducible = 1 equation}) is greater than 1. Hence, to test whether an arbitrary representation is reducible, calculate the RHS of Eq. (\ref{irreducible = 1 equation}) and see whether it is greater than or equal to 1; the former situation suggests that the representation is reducible, whereas the latter irreducible. 
	
	
	\subsection{Special Unitary Group $SU(N)$}
	Lie groups are groups that have infinitely many group elements that are also differentiable manifolds, which allows us to do differential analysis (e.g., infinitesimal transformation). A Lie group is special unitary if it satisfies the following:
	\begin{align}
		U^\dagger U = 1 \label{unitary of su n}; \\
		\det(U) = 1 \label{det 1 of su n}.
	\end{align}
	
	In the fundamental representation of $SU(N)$ group, each of the indices $a$, $b$ and $c$ in tensor $\psi^{ab}_{c}$ runs from $1$ to $N$. For example, a rank-$1$ tensor $\psi^a$ can be expressed by a column vector:
	\begin{equation}
		\psi^a = \begin{pmatrix}
			\psi^1 \\
			\psi^2 \\
			\vdots\\
			\psi^N
		\end{pmatrix}.
	\end{equation}
	Different from that of $SO(N)$ groups that feature real orthogonal matrices, the quadratic invariant in $SU(N)$ reads
	\begin{equation}
		\zeta^\dagger\psi = \sum_{j=1}^N\zeta^{j*}\psi^j \equiv \zeta_j\psi^j,
	\end{equation}
	where $\zeta_j$ with lower indices is called covariant tensor and equal to complex conjugate of contravariant tensor $\zeta^j$. A covariant tensor transforms as
	\begin{equation}
		\zeta_i \rightarrow \zeta_j(U^\dagger)^{j}_{\  i}.
	\end{equation}
	Thus, in $SU(N)$, all other operations also need to abide the convention of upper lower indices. For example, taking trace:
	\begin{equation}
		\psi^{ij}_j \rightarrow U^i_{\ l}U^{j}_{\ m}(U^\dagger)^n_{\ j}\psi^{lm}_n = U^i_{\ l}\psi^{lm}_m.
	\end{equation}
	
	
	
	
	
	Any Lie group element that is continuously connected to the identity can be written in terms of the identity multiplied by the exponential of generator $T^a$, which expresses the infinitesimal transformation under the given representation:
	\begin{equation}
		U = \exp(i\theta^aT^a) \cdot I.
	\end{equation}
	The algebra of these generators is called a Lie algebra. Lie algebra is defined by commutation relations of the generators:
	\begin{equation}
		\left[ T^a , T^b \right] = if^{abc}T^c,
	\end{equation}
	where $[A,B] \equiv AB-BA$, which measures the commutativity of the given transformation. $f^{abc}$ is called the structure constant. Since $[T^a, T^b] \propto f^{abc} = -[T^b, T^a] \propto f^{bac}$ in $SU(N)$, $f^{abc}$ is anti-symmetric:
	\begin{equation}
		f^{abc} = -f^{bac}.
	\end{equation}
	The Jacobi identity:
	\begin{equation}
		[A, [B, C]] + [B, [C, A]] + [C, [A, B]] = 0 \Rightarrow f^{abd}f^{dce} + f^{bcd}f^{dae} + f^{cad}f^{dbe} = 0.
	\end{equation}
	
	
	
	Substitute $U = \exp(i\theta^aT^a)$ into Eq. (\ref{unitary of su n}), we get 
	\begin{equation}
		T^\dagger = T.
	\end{equation}
	In other words, unitary condition enforces hermiticity of the generator. Moreover, by utilizing the fact that $T$ can be diagonalized and that $\det(U) = 1$, we can derive (\cite{azee2016group}, p. 236)
	\begin{equation}
		\tr(H) = 0.
	\end{equation}
	Therefore, we know that a generator of $SU(N)$ needs to be \textbf{hermitian} and \textbf{traceless}. 
	
	To calculate free elements in a $SU(N)$ generator:
	\begin{enumerate}
		\item there are $N-1$ free real indices on the diagonal because of the traceless and hermitian conditions;
		\item for the $N(N-1)$ matrix elements lefted, hermiticity enforces the number of free (complex) indices to be $\frac{1}{2}N(N-1)$, which means that $N(N-1)$ real numbers are needed to define the matrix. 
	\end{enumerate}
	In total, $2 \cdot \frac{1}{2}N(N-1)+N-1 = N^2-1$ real numbers are required to define an $N\times N$ traceless hermitian matrix. And, since we need $N^2-1$ free parameters to specify such a matrix, all matrices in the space spanned by traceless hermitian matrices can be written as a linear combination of $N^2 - 1$ linearly independent basis matrices. The dimensionality of the space is thus also $N^2 - 1$. 
	
	Define the adjoint representation of group $SU(N)$ as the representation of  $SU(N)$ group elements that act on the Lie algebra of the $SU(N)$. That is, we want to see how $SU(N)$ group elements (the complex "rotation") would act on the generators (for example, $T^a$). In $SU(N)$, each of the $N^2 - 1$ generators is a rank-2 traceless hermitian tensor (i.e., a tracless hermitian matrix). Therefore, let $\psi^i_j$ be a tensor in the space spanned by the $N^2 - 1$ generators  (i.e., $\psi^i_j$ need not to be any of the generators themselves, but could be their linear combinations).According to the rule by which a tensor transforms,
	\begin{equation}
		\psi'^k_{\ l} = U^k_{\ i}\psi^i_{\ j} (U^\dagger)^j_{\ l} = U\psi U^\dagger,
	\end{equation}
	in which $U$ is the fundamental representation of $SU(N)$, and the transformation corresponds to the adjoint action $\text{Ad}_U(\psi) = U\psi U^\dagger$. To be able to act on $N^2 - 1$ dimensional tensor $\psi$, the dimensionality of the adjoint representation of $SU (N)$ is also $N^2 -1$. 
	
	Choose a basis $\{T_a\}_{a=1}^{N^2-1}$ of $\psi$ with
	\begin{align}
		\tr(T_aT_b)=\frac12\delta_{ab}, \\ [T_a,T_b]=i f_{abc}T_c,
	\end{align}
	where $f_{abc}$ are the real, totally antisymmetric structure constants, and $\frac{1}{2}$ being normalization factor. Then the matrix elements of $\text{Ad}_U$ in this basis are
	\begin{equation}
		(\text{Ad}_U)_{ab}=2\,\tr(T_a U T_b U^\dagger),
	\end{equation}
	so that $\text{Ad}_U(T_b)=\sum_a (\text{Ad}_U)_{ab}T_a$. 
	
	For an infinitesimal transformation $U=e^{i\theta_c T_c}\approx 1+i\theta_c T_c$, we obtain
	\[
	\text{Ad}_U(T_b)\approx T_b + i\theta_c[T_c,T_b]=T_b-\theta_c f_{cbd}T_d,
	\]
	hence the Lie algebra generators in the adjoint representation are given by
	\[
	(\text{Ad}_{T^a})_{bc}=-i f_{abc},
	\]
	satisfying $[\text{Ad}_{T^a},\text{Ad}_{T^b}]=i f_{abc}\text{Ad}_{T^c}$. Therefore the adjoint representation is a faithful representation of $\mathfrak{su}(N)$ on itself.
	
	Some identities in $SU(N)$ (\cite{schwartz1}, p. 487, 488):
	\begin{align}
		&\tr(T^aT^b) = T^a_{ji}T^b_{ij} = T_F\delta^{ab}; \nonumber\\
		&\sum_a(T^aT^a)_{ij} = C_F\delta_{ij};\nonumber\\
		&f^{acd}f^{bcd} = C_A\delta^{ab};\nonumber\\
		&f^{abc} = -\frac{i}{T_F}\tr([T^a, T^b]T^c);\nonumber \\
		&\sum_aT^a_{ij}T^a_{kl} = \frac{1}{2}\left( \delta_{il}\delta_{kj}-\frac{1}{N}\delta_{ij}\delta_{kl}\right) \text{ (Fierz Identity)},\label{fabc formula & Fierz Identity}
	\end{align}
	where $T_F = \frac{1}{2}$ is called the index of the representation, $C_F = \frac{N^2-1}{2N}$ is the quadratic Casimir for the fundamental representation of $SU(N)$, and $C_A = N$ is the quadratic Casimir for the adjoint representation of $SU(N)$. 
	
	
	
	%===============
	
	\subsection{$SU(3)$ Symmetry of Color}
	The color plays the role of charge in quantum chromodynamics. A quark can be red, green, or blue. An anti-quark can be anti-red, anti-green, or anti-blue. In QCD, the strong interaction follows the color $SU(3)$ symmetry--the strong interaction is invariant under rotations in color space. This $SU(3)$ symmetry is exact.
	
	The fundamental representation of $SU(3)$ is three-dimensional. It corresponds to the quark with three colors: red, blue, and green. The generators of the fundamental representation are the eight $3\times 3$ traceless Hermitian matrices $T^a = \frac{1}{2}\lambda^a$, with normalization $\tr(T^aT^b) = \frac12 \delta^{ab}$ (\cite{schwartz1}, p. 485):
	\begin{align}
		\lambda^1 &= \begin{pmatrix} 0 & 1 & 0 \\ 1 & 0 & 0 \\ 0 & 0 & 0 \end{pmatrix}, \quad
		\lambda^2 = \begin{pmatrix} 0 & -i & 0 \\ i & 0 & 0 \\ 0 & 0 & 0 \end{pmatrix}, \quad
		\lambda^3 = \begin{pmatrix} 1 & 0 & 0 \\ 0 & -1 & 0 \\ 0 & 0 & 0 \end{pmatrix}, \nonumber\\
		\lambda^4 &= \begin{pmatrix} 0 & 0 & 1 \\ 0 & 0 & 0 \\ 1 & 0 & 0 \end{pmatrix}, \quad
		\lambda^5 = \begin{pmatrix} 0 & 0 & -i \\ 0 & 0 & 0 \\ i & 0 & 0 \end{pmatrix}, \nonumber\\
		\lambda^6 &= \begin{pmatrix} 0 & 0 & 0 \\ 0 & 0 & 1 \\ 0 & 1 & 0 \end{pmatrix}, \quad
		\lambda^7 = \begin{pmatrix} 0 & 0 & 0 \\ 0 & 0 & -i \\ 0 & i & 0 \end{pmatrix}, \quad
		\lambda^8 = \frac{1}{\sqrt{3}}\begin{pmatrix} 1 & 0 & 0 \\ 0 & 1 & 0 \\ 0 & 0 & -2 \end{pmatrix}. \label{Gellmann mtx}
	\end{align}
	Any group element of $SU(3)$ in the fundamental representation can be written as
	\begin{equation}
		U = \exp(i\theta^aT^a) = \exp(i\theta^a\lambda^a/2).
	\end{equation}
	For example, $U = e^{i\theta T^3} = \begin{pmatrix}
		e^{-i\theta/2} &  & \\
		& e^{i\theta/2} & \\
		& & 1
	\end{pmatrix}.$
	The quadratic Casimir of the fundamental representation is
	\begin{equation}
		C_F = \frac{3 \times 3 - 1}{2 \times 3} = \frac{4}{3}.
	\end{equation}
	
	A gluon takes a combination of color and anti-color, for example $r\bar{r}$. That is, a gluon is described by the direct product of the 3-dimensional fundamental representation $\mathbf{3}$ and its conjugate $\bar{\mathbf{3}}$, i.e.\ $\mathbf{3} \otimes \bar{\mathbf{3}}$. Each index transforms independently, hence a gluon field is mathematically a rank-2 tensor $(G_\mu)^i_j$, with the upper index (contravariant) representing the color and the lower index (covariant) representing the anti-color. The 9-dimensional product space decomposes into two irreducible representations:
	\begin{equation}
		\mathbf{3} \otimes \bar{\mathbf{3}} = \mathbf{8} \oplus \mathbf{1},
	\end{equation}
	physically corresponding to the traceless part (the octet, containing the eight gluons) and the trace part (a color singlet that does not mediate strong interactions).
	
	The octet forms the adjoint representation of $SU(3)$, whose dimension is $3^2-1 = 8$. The matrix elements of the generators in the adjoint representation are
	\begin{equation}
		(T^a_{\text{adj}})_{bc} = -if^{abc},
	\end{equation}
	and the quadratic Casimir is $C_A = 3$. The adjoint representation acts on the Lie algebra itself, which means the gluon fields transform as an octet under $SU(3)$. This property describes the self-interaction of gluons, a unique feature of non-Abelian gauge theories like QCD.
	
	\subsection{Structure Constant}
	The structure constants $f^{abc}$ are defined by the commutation relations $[T^a, T^b] = i f^{abc} T^c$. They can be computed from the trace formula that is already introduced in Eq. (\ref{fabc formula}),
	\begin{equation}
		f^{abc} = -\frac{i}{T_F}\,\tr\!\big([T^a, T^b] T^c\big),
	\end{equation}
	with $T_F = \frac12$ in the fundamental representation. For example, to find $f^{123}$ we evaluate the commutator of $T^1$ and $T^2$:
	\[
	[T^1, T^2] = \frac{1}{4}[\lambda^1, \lambda^2] = \frac{1}{4}(2i\lambda^3) = i T^3.
	\]
	Inserting this into the trace formula gives
	\[
	f^{123} = -\frac{i}{1/2}\,\tr\!\big(i T^3 T^3\big) = 2\,\tr(T^3 T^3) = 2\cdot\frac12 = 1.
	\]
	Proceeding similarly one obtains all non-zero structure constants (up to permutations), for example
	\[
	f^{147} = f^{246} = f^{257} = f^{345} = \frac12,\qquad
	f^{156} = f^{367} = -\frac12,\qquad
	f^{458} = f^{678} = \frac{\sqrt{3}}{2}.
	\]
	The normalization condition for the structure constants in the adjoint representation reads
	\begin{equation}
		f^{abc}f^{dbc} = C_A \delta^{ad} = 3\delta^{ad}.
	\end{equation}
	
	\subsection{Color Wave Function}
	The wave function in quantum mechanics represents a quantum state. For a hadron, its wave function can be decomposed into the direct product of spatial, flavor, spin, and color wave function: 
	\begin{equation}
		\Psi_{\text{total}} = \Psi_{\text{color}} \otimes \Psi_{\text{spin}}\otimes \Psi_{\text{flavor}} \otimes \Psi_{\text{space}}.
	\end{equation}
	A quark carries a color, whether $\ket{r}$, $\ket{g}$, $\ket{b}$, or their linear combinations. An anti-quark carries an anti-color, whether $\ket{\bar r}$, $\ket{\bar g}$, $\ket{\bar b}$, or their linear combinations. A gluon carries a combination of color and anticolor, such as the state
	\begin{equation}
		\ket{1} = \frac{1}{\sqrt{2}}(\ket{r\bar b} + \ket{b \bar r}).
	\end{equation}
	As mentioned earlier, the $3 \times 3 = 9$ dimensional space for gluons can be decomposed into an $8$ dimensional space (traceless), with basis vector $\ket{1}, \ket{2},...,\ket{8}$ and a scalar (with trace) (\cite{griffiths1987IEP}, p. 280)
	\begin{equation}
		\ket{9} = \frac{1}{\sqrt{3}}\left( \ket{r\bar r} + \ket{g \bar g} + \ket{b \bar b}\right). \label{singlet config}
	\end{equation} Color confinement requires all non-singlet state unseparable; together with the experimental fact that strong interaction is short-range interaction and that we never detected a single gluon, one may conclude that a singlet gluon does not exist. Since a singlet state can only emit or absorb singlet states, interactions between baryons or mesons (which are all singlets to be observable) would be long-range forces if a singlet gluon existed. Hence, as a bound state of multiple gluons, a glueball (which we are currently investigating) is a color singlet (\cite{griffiths1987IEP}, p. 280, 281; \cite{glueballhunting}, p.2).
	
	
	Define the Casimir operator to be:
	\begin{equation}
		\hat{C}_R = T^a_R T^a_R.
	\end{equation}
	It can be proved that the Casimir operator commutes to all generators (\cite{schwartz1}, p. 486):
	\begin{equation}
		[\hat{C}_R, T^b_R] = 0.
	\end{equation}
	Therefore, by Schur's lemma (see section \ref{group theory basics section}), the Casimir is proportional to the identity:
	\begin{equation}
		\hat{C}_R = C_2(R)\cdot\mathbb{I},
	\end{equation}
	where $C_2(R)$ is a constant named quadratic Casimir. 
	
	Note that the quadratic Casimir is of the representation. In the $SU(3)$ fundamental representation $\mathbf{3}$ and $\mathbf{\bar 3}$, the quadratic Casimir
	\begin{equation}
		C_F = \frac{N^2-1}{2N} = \frac{4}{3};
	\end{equation}
	in the $SU(3)$ adjoint representation, the quadratic Casimir
	\begin{equation}
		C_A = N = 3. 
	\end{equation}
	
	
	
	Especially, the quadratic Casimir $C_\mathbf{1} = 0$ for singlet representation. This is because the color singlet corresponds to the traceless $\mathbf{1}$ representation of $SU(3)$, 
	and, by definition, it is invariant under any $SU(3)$ group transformation:
	\begin{equation}
		T^a |\mathbf{1}\rangle = 0, \text{ with } a = 1,2,...,8.
	\end{equation}
	Because each individual generator gives zero when acting on the singlet, the Casimir also gives zero:
	\begin{equation}
		\hat{C}|\mathbf{1}\rangle 
		= T^a \left(T^a |\mathbf{1}\rangle\right)  
		= 0.
	\end{equation}
	In contrast, for any non-singlet representation $\mathbf{R}$, the Casimir acts as a positive constant $C_2(\mathbf{R}) > 0$ times the state. Therefore, we use Casimir to test whether a state is a color singlet (\cite{Lijun1}, p. 29).
	
	\subsubsection{Gluon Singlet}
	A gluon carries a color and an anti-color, for example $\ket{r\bar b}$; the gluon states live in a $3 \times 3 = 9$ dimensional space (we do not separate the singlet state yet because it is what we're looking for):
	\begin{align}
		\ket{g_1} &= \ket{r\bar r}, &
		\ket{g_2} &= \ket{r\bar g}, &
		\ket{g_3} &= \ket{r\bar b}, \\
		\ket{g_4} &= \ket{g\bar r}, &
		\ket{g_5} &= \ket{g\bar g}, &
		\ket{g_6} &= \ket{g\bar b}, \\
		\ket{g_7} &= \ket{b\bar r}, &
		\ket{g_8} &= \ket{b\bar g}, &
		\ket{g_9} &= \ket{b\bar b}.
	\end{align}
	To describe how gluon states transform, we need to work in the adjoint representation of $SU(3)$, the nine-dimensional generators $T^a_{\mathbf{3} \otimes \bar{\mathbf{3}}}$ of which can be constructed by generators in the fundamental representation $T^a$ as
	\begin{align}
		T^a_{\mathbf{3} \otimes \bar{\mathbf{3}}} = T^a \otimes \mathbb{I}_3 + \mathbb{I}_3 \otimes \tilde{T}^a \\= T^a \otimes \mathbb{I}_3 + \mathbb{I}_3 \otimes (-T^a)^T.
	\end{align}
	Since the operator on the left of $\otimes$ acting only on $\ket{i}$ in $\ket{i, \bar j} = \ket{i}\otimes\ket{j}$ and the operator on the right acting only on $\ket{\bar j}$, we have
	\begin{align}
		(T^a_{\mathbf{3} \otimes \bar{\mathbf{3}}})_{(i, j), (k, l)} = (T^a)_{ik} \otimes (\mathbb{I}_3)_{j l} + (\mathbb{I}_3)_{ik} \otimes (-T^a)_{l j} = (T^a)_{ik} \delta_{j l} - \delta_{ik}(T^a)_{lj}
	\end{align}
	Therefore, with the help of Fierz identity in Eq. (\ref{fabc formula & Fierz Identity}), the Casimir 
	\begin{align*}
		(\hat{C})_{(i, j), (m, n)} &= (T^a_{\mathbf{3} \otimes \bar{\mathbf{3}}})_{(i, j), (k, l)}(T^a_{\mathbf{3} \otimes \bar{\mathbf{3}}})_{(k, l), (m, n)}\\\
		&= \bigl[(T^a)_{ik} \delta_{jl} - \delta_{ik}(T^a)_{lj}\bigr]
		\bigl[(T^a)_{km} \delta_{ln} - \delta_{km}(T^a)_{nl}\bigr] \\
		&= (T^a)_{ik}(T^a)_{km} \delta_{jl}\delta_{ln}
		- (T^a)_{ik}(T^a)_{nl} \delta_{jl}\delta_{km} \\
		&\quad - (T^a)_{lj}(T^a)_{km} \delta_{ik}\delta_{ln}
		+ (T^a)_{lj}(T^a)_{nl} \delta_{ik}\delta_{km} \\
		&= (T^a T^a)_{im} \delta_{jn} - (T^a)_{im}(T^a)_{nj}
		- (T^a)_{im}(T^a)_{nj} + \delta_{im}(T^a T^a)_{nj} \\
		&= C_F \delta_{im}\delta_{jn} - 2(T^a)_{im}(T^a)_{nj} + C_F \delta_{im}\delta_{jn} \\
		&= 2C_F \delta_{im}\delta_{jn} - 2(T^a)_{im}(T^a)_{nj} \\
		&= 2\cdot\frac{4}{3}\,\delta_{im}\delta_{jn}
		- 2\cdot\frac12\Bigl( \delta_{ij}\delta_{mn} - \frac13\delta_{im}\delta_{jn} \Bigr) \\
		&= \frac{8}{3}\delta_{im}\delta_{jn} - \delta_{ij}\delta_{mn} + \frac13\delta_{im}\delta_{jn} \\
		&= 3\delta_{im}\delta_{jn} - \delta_{ij}\delta_{mn},
	\end{align*}
	or equivalently,
	\begin{align}
		(\hat{C})_{(i, j), (m, n)} = \begin{pmatrix}
			2&0&0&0&-1&0&0&0&-1\\
			0&3&0&0&0&0&0&0&0\\
			0&0&3&0&0&0&0&0&0\\
			0&0&0&3&0&0&0&0&0\\
			-1&0&0&0&2&0&0&0&-1\\
			0&0&0&0&0&3&0&0&0\\
			0&0&0&0&0&0&3&0&0\\
			0&0&0&0&0&0&0&3&0\\
			-1&0&0&0&-1&0&0&0&2
		\end{pmatrix}.
	\end{align}
	From previous conclusions, we know that a Casimir acting on a color singlet would result in zero. Therefore, to find the expression of the singlet gluon, we need to diagonalize $\hat{C}$ and select the eigenstates (which are $9 \times 1$ column matrices) corresponding to the eigenvalue $0$. The result shows that there is only one eigenvalue $\text{(eigen value)}_9 = 0$, with the eigenvector
	\begin{equation}
		\text{(eigen vector)}_9 = \begin{pmatrix}
			1 \\ 0\\ 0\\ 0 \\ 1 \\ 0 \\ 0 \\ 0 \\ 1
		\end{pmatrix} = \frac{1}{\sqrt{3}}\left( \ket{r\bar r} + \ket{g \bar g} + \ket{b \bar b}\right),
	\end{equation}
	where $\frac{1}{\sqrt{3}}$ is the normalization factor. This result perfectly matches Eq. (\ref{singlet config}).
	
	
	
	\subsection{Cartan Sub-Algebra and Cartan-Weyl Basis}
	
	The eight Gell-mann matrices, as listed in Eq. (\ref{Gellmann mtx}), were the generators of $SU(3)$ fundamental representation with appropriate factors. Among the eight Gell-mann matrices, there are two diagonal matrices: $\lambda^3$ and $\lambda^8$. Obviously, they commute with each other:
	\begin{equation}
		[\lambda^3, \lambda^8] = 0 \Rightarrow [T^3, T^8] = 0. 
	\end{equation}
	
	The Cartan sub-algebra of $SU(3)$ is defined as the algebra spanned by $T^3$ and $T^8$. From quantum mechanics we know that there exists a complete set of simultaneous eigenfunctions for two compatible operators. Therefore, it is legitimate to introduce quantum numbers $i_3$ and $Y$ as eigenvalues of $T^3$ and $T^8$, respectively:
	\begin{equation}
		\psi = |i_3, Y\rangle. 
	\end{equation}
	It is illuminating to compare the two quantum numbers to the $z$-component of angular momentum in quantum mechancis $j_z = m$. The only difference is that the visuallization of $m$ is a ladder, whereas $(i_3, Y)$, being two-dimensional, lies on a plane.
	
	A smart way to deal with these generators requires us to replace the Gell-mann matrices by operators under the Cartan-Weyl basis. For a general case of $n$ generators $X^n$, we first pick up \textbf{as much as possible} mutually compatible generators (suppose there are $l$ such generators):
	\begin{equation}
		\left[H^i, H^j\right] = 0,
	\end{equation}
	where we call $l$ the rank of this representation. The lefted $n - l$ generators will be rearranged so that
	\begin{equation}
		\left[H^i, E_{\vec\alpha} \right] = \alpha^i E_{\vec\alpha},
	\end{equation}
	which means that the $E_{\vec{\alpha}}$ acts as a ladder operator; if $H\ket{\mu} = \mu \ket{\mu}$, then $HE\ket{\mu} = [H, E]\ket{\mu} + EH\ket{\mu} = (\alpha^i + \mu) E\ket{\mu} \propto H\ket{\mu + \alpha}$. $\vec{\alpha} = (\alpha^1, ..., \alpha^i$,...) is called the root vector, which characterizes how the ``ladder'' operators move the states. 
	
	In $SU(3)$ fundamental representation, the $H$'s are $T^3$ and $T^8$, and the $E$'s are (\cite{azee2016group}, p. 328)
	\begin{align}
		&I_{\pm} = \frac{1}{2}\left( \lambda^1 \pm i\lambda^2 \right) \rightarrow \text{root vector: }(\pm 1, 0); \\
		&U_{\pm} = \frac{1}{2}\left( \lambda^6 \pm i\lambda^7 \right) \rightarrow \text{root vector: }\left(\mp \frac{1}{2},\ \pm \frac{\sqrt{3}}{2}\right); \\
		&V_{\pm} = \frac{1}{2}\left( \lambda^4 \pm i\lambda^5 \right) \rightarrow \text{root vector: }\left(\pm \frac{1}{2},\ \pm \frac{\sqrt{3}}{2}\right).
	\end{align}
	
	
	
	
	\bibliographystyle{unsrt}
	\bibliography{main}
\end{document}